\section{Consequences and limits}\label{sec:discussion}

The main reduction identifies the obstacle to an exact separation
oracle.  It is not positive semidefiniteness testing, triangle
separation, or recognition of a small fixed family of inequalities.
Every reduction point passes those tests.  Its remaining violation
encodes the choice of a cover, and the proof excludes all other
integer coefficient vectors.  Polynomial-time enforcement of the
standard relaxations therefore leaves an NP-complete separation
problem.

This distinction matters when a solver adds only selected members of
an infinite class.  Any fixed gonality gives a polynomially enumerable
family, but the exponent grows with that gonality.  Theorem
\ref{thm:relaxation-hardness} proves that increasing a fixed cutoff
does not make the unrestricted separation problem tractable.  The
rank algorithm offers a different restriction: it exploits the
dimension of the rational lattice induced by the input matrix.  The
threshold algorithm restricts the norm of a potentially violated
split normal by asking for a sufficiently large normalized violation.
These are useful and distinct ways to limit a search.  Neither
provides a polynomial-time exact oracle at unrestricted rank and
dimension.

The numerical scale of the construction is also explicit.  On a yes
instance, every base hypermetric or rounded positive semidefinite
violation has cut-scale value $1/(2N)$, and every corresponding
divided Boolean quadric cut or clique violation has value $1/(4N)$.
The orthogonal extension has pure and positive odd clique violations
of the same inverse-polynomial order.  All coordinates have a
polynomial common denominator.  Hardness therefore does not depend
on exponentially small slacks.  The exact zero-or-positive maximum
also rules out finite-factor approximation of raw violation and
additive approximation at the instance-dependent accuracy in
\cref{cor:raw-approximation}.  It does not prove hardness at a fixed
additive tolerance or show that every hard instance must have small
violations.

Gap-$0$ separation concerns points outside the elliptope: its
inequalities have zero right-hand side in the $Z$ representation and
hold whenever $Z\succeq0$.  The remaining classification problem is
full gap separation at positive semidefinite points.  The present
reductions control violations with gap $1$ but do not exclude a
violator with gap at least $3$ on a no-instance.  A reduction for the
full family must therefore control this larger right-hand side.
